\chapter{Marketing Information Systems and the Sales Order Process}
\label{ch:sales-order}

\chapterguide
  {You should understand functional integration and the idea of shared transaction data.}
  {describe an end-to-end sales order process, diagnose problems caused by disconnected systems, explain the order-to-cash sequence, and show how sales, inventory, and finance are linked.}

\begin{sourcebasis}
This chapter uses Tutorial 5 and Chapter 3 of \emph{Concepts in Enterprise
Resource Planning}. The Fitter Snacker case is retained as the main example of
an inefficient, unintegrated sales process; the Odoo sections map the same
concepts onto the West End Bicycles practical.
\end{sourcebasis}

\section{What the sales order process includes}

The \term{sales order process} runs from a customer's request through fulfilment
and payment. It is much broader than entering an order: quotation, order
capture, credit terms, availability checking, picking and delivery, invoicing,
payment, and sometimes returns.

Each step is owned by a different function and each produces information the
next step depends on.

\begin{erpfig}[Order-to-cash by owner]{The sales order process with the owner of
each step and the record it creates. Reading down the third row shows why no
single department can be held responsible for the customer's
experience.}{fig:order-to-cash-lanes}
\begin{tikzpicture}[x=1cm,y=1cm,
  tight/.style={-{Stealth[length=1.6mm]},semithick,draw=UTPTeal}]
\node[erpcap,anchor=east,align=right] at (1.05,2.15) {\bfseries Step};
\node[erpcap,anchor=east,align=right] at (1.05,1.12) {\bfseries Owner};
\node[erpcap,anchor=east,align=right] at (1.05,0.12) {\bfseries Creates};
\foreach \i/\s/\o/\r/\c in {%
  0/Quotation/{Sales}/{offer, price}/UTPBlue,
  1/{Sales order}/{Sales}/{commitment}/UTPBlue,
  2/{Availability, credit}/{Inventory, Finance}/{reservation}/UTPTeal,
  3/{Pick, pack, deliver}/{Inventory}/{shipped qty}/UTPTeal,
  4/Invoice/{Finance}/{receivable}/UTPGreen,
  5/Payment/{Finance}/{settlement}/UTPGreen}{
  \pgfmathsetmacro{\x}{2.2+\i*2.35}
  \node[erpstep,minimum width=2.0cm,minimum height=0.86cm,draw=\c,
        fill=\c!12,text width=1.75cm] (s\i) at (\x,2.15) {\s};
  \node[erpdoc,minimum width=2.0cm,minimum height=0.66cm,draw=\c!60,text width=1.75cm] (o\i) at (\x,1.12) {\o};
  \node[erpchip,text width=1.6cm,align=center] (r\i) at (\x,0.12) {\r};
  \draw[erplink] (\x,1.72) -- (\x,1.45);
  \draw[erplink] (\x,0.79) -- (\x,0.42);
}
\foreach \i in {0,...,4}{\pgfmathtruncatemacro{\j}{\i+1}\draw[tight] (s\i)--(s\j);}
\draw[erpback] (s5.north) to[out=105,in=75]
  node[midway,above,erpnote]{payment status returns to Sales} (s0.north);
\end{tikzpicture}
\end{erpfig}

\section{Fitter Snacker: what goes wrong without integration}

Fitter Snacker runs separate sales, warehouse, and accounting systems. Data
crosses between them by fax, paper report, and periodic file transfer. Every
recurring problem traces back to one structural fact: \textbf{each department
holds its own copy of the order, and the copies drift apart.}

\begin{erpfig}[Copies versus one record]{The same three departments, two ways.
On the left, each holds a copy and the copies disagree between transfers; on the
right, one record serves all three.}{fig:fitter-snacker-silos}
\begin{tikzpicture}[x=1cm,y=1cm]
% ---- left: siloed ----
\node[erphead,color=UTPRed] at (2.6,3.3) {Separate systems};
\node[erpbad,minimum width=2.6cm,minimum height=0.95cm] (s1) at (2.6,2.2)
  {Sales\\\tiny\mdseries own order copy};
\node[erpbad,minimum width=2.6cm,minimum height=0.95cm] (s2) at (2.6,0.75)
  {Warehouse\\\tiny\mdseries own stock copy};
\node[erpbad,minimum width=2.6cm,minimum height=0.95cm] (s3) at (2.6,-0.7)
  {Accounting\\\tiny\mdseries own invoice copy};
\draw[erpslow] (s1) -- node[erpnote,right=1mm,text width=2.5cm,align=left]
  {fax, batch transfer\\ \emph{hours or days}} (s2);
\draw[erpslow] (s2) -- node[erpnote,right=1mm,text width=2.5cm,align=left]
  {re-keyed by hand\\ \emph{may not match}} (s3);
% ---- right: integrated ----
\node[erphead,color=UTPGreen] at (11.6,3.3) {One shared record};
\node[erpstore,draw=UTPBlue,minimum width=2.6cm,minimum height=1.35cm] (h) at (11.6,0.75)
  {Shared\\\scriptsize\bfseries transaction\\\scriptsize\bfseries data};
\node[erpgood,minimum width=2.3cm,minimum height=0.75cm] (g1) at (8.4,2.2) {Sales};
\node[erpgood,minimum width=2.3cm,minimum height=0.75cm] (g2) at (8.4,-0.7) {Warehouse};
\node[erpgood,minimum width=2.3cm,minimum height=0.75cm] (g3) at (14.7,0.75) {Accounting};
\draw[erpflow] (g1) -- (h); \draw[erpflow] (g2) -- (h); \draw[erpflow] (h) -- (g3);
\node[erpnote,text width=4.6cm] at (11.6,-1.35)
  {every function reads and writes the \emph{same} order};
\end{tikzpicture}
\end{erpfig}

Read the five textbook symptoms as one cause seen from five desks.

\begin{erptable}
\begin{tabularx}{\linewidth}{@{}>{\bfseries\leavevmode\color{ERPNavy}}p{0.20\linewidth}X>{\itshape}p{0.26\linewidth}@{}}
\toprule
\rowcolor{ERPPaleBlue!92}
Symptom & What actually happens & The stale copy \\
\midrule
Quotation and pricing & Arithmetic slips, discounts combined wrongly, terms the office has not seen; faxes misread into duplicate orders. & price \\
Availability and dates & The clerk asks the warehouse for a ship date built on stock figures that are already out of date. & stock \\
Credit & The customer is judged against a periodic printed report, so a recent payment is invisible --- good customers refused, over-exposed ones approved. & receivables \\
Order filling & Paper picking lists travel in batches; stock moves nobody writes down never reach the records. & inventory \\
Invoicing & Order data transfers to accounting on a schedule, so fulfilment changes after the transfer are billed wrongly. & shipped quantity \\
\bottomrule
\end{tabularx}
\end{erptable}

The cost shows up as elapsed time. Most of an unintegrated order's life is spent
waiting for a transfer, not being worked on.

\begin{erpfig}[Where the days go]{The same order under both processes. The
shaded blocks are work; the gaps are waiting for a batch, a fax, or a report.
Integration removes the gaps rather than making anyone type
faster.}{fig:fitter-timeline}
\begin{tikzpicture}[x=1cm,y=1cm]
\draw[draw=UTPMuted!45] (0,-1.55) -- (14.6,-1.55);
\foreach \d in {0,1,...,7}{
  \pgfmathsetmacro{\x}{\d*1.82}
  \draw[draw=UTPMuted!45] (\x,-1.62)--(\x,-1.48);
  \node[erpnote,below=1mm] at (\x,-1.55) {day \d};}
% old process
\node[erpcap,anchor=west,color=UTPRed] at (0,1.15) {\bfseries Fitter Snacker (unintegrated)};
\foreach \a/\b/\t in {0/1.1/Quote, 3.3/4.4/{Order entered}, 6.6/8.1/{Credit checked},
                      9.6/11.0/{Picked}, 12.4/14.2/{Invoiced}}{
  \draw[draw=UTPRed,fill=UTPPaleRed,rounded corners=1mm] (\a,0.05) rectangle (\b,0.75);
  \node[erpnote,anchor=center] at ({(\a+\b)/2},0.4) {\t};}
\foreach \a/\b in {1.1/3.3, 4.4/6.6, 8.1/9.6, 11.0/12.4}{
  \draw[draw=UTPRed!55,dashed] (\a,0.4)--(\b,0.4);
  \node[erpnote,color=UTPRed,above=0.5mm] at ({(\a+\b)/2},0.4) {wait};}
% integrated
\node[erpcap,anchor=west,color=UTPGreen] at (0,-0.05) {\bfseries Integrated (one shared record)};
\foreach \a/\b/\t in {0/1.1/Quote, 1.1/2.1/{Order}, 2.1/3.3/{Credit},
                      3.3/4.5/{Picked}, 4.5/5.7/{Invoiced}}{
  \draw[draw=UTPGreen,fill=UTPPaleTeal,rounded corners=1mm] (\a,-1.15) rectangle (\b,-0.45);
  \node[erpnote,anchor=center] at ({(\a+\b)/2},-0.8) {\t};}
\node[erpchip,anchor=west] at (6.1,-0.8) {same work, no transfers to wait for};
\end{tikzpicture}
\end{erpfig}

\begin{pitfall}
None of this is careless staff; it is process design. A process that requires
people to copy, fax, re-key, and reconcile the same fact across systems produces
errors and delays no matter how conscientious they are. Fix the structure.
\end{pitfall}

\section{What an integrated ERP changes}

An integrated system replaces separate copies with a linked document flow. The
value is not that the documents exist --- Fitter Snacker had documents too ---
but that each one inherits its data from the one before it.

\begin{erpfig}[Document lineage]{What each document carries forward. Nothing in
the lower row is re-typed; it is inherited. The return arrow is the answer Sales
can give a customer without telephoning another department.}{fig:document-lineage}
\begin{tikzpicture}[x=1cm,y=1cm]
\foreach \i/\d/\c in {0/Quotation/UTPBlue, 1/{Sales Order}/UTPBlue,
                      2/Delivery/UTPTeal, 3/Invoice/UTPGreen, 4/Payment/UTPGreen}{
  \pgfmathsetmacro{\x}{\i*3.15}
  \node[erpstep,minimum width=2.75cm,minimum height=0.8cm,draw=\c,fill=\c!12] (d\i) at (\x,1.5) {\d};}
\foreach \i in {0,...,3}{\pgfmathtruncatemacro{\j}{\i+1}\draw[erpflow] (d\i)--(d\j);}
\foreach \i/\t in {0/{customer\\ product\\ quantity\\ unit price},
                   1/{$+$ confirmed\\ commitment\\ $+$ salesperson},
                   2/{$+$ what actually\\ left the\\ warehouse},
                   3/{$+$ amount due\\ $+$ terms\\ $+$ tax},
                   4/{$+$ cash received\\ $+$ settlement\\ date}}{
  \pgfmathsetmacro{\x}{\i*3.15}
  \node[erpdoc,minimum width=2.75cm,text width=2.4cm,align=center,minimum height=1.5cm] (c\i) at (\x,-0.35) {\t};
  \draw[erplink] (\x,1.1) -- (\x,0.42);}
\node[erpcap,color=UTPGreen] at (6.3,-1.55) {\bfseries what each document carries forward};
\draw[erpback] (d4.north) to[out=120,in=60] node[midway,above,erpnote]
  {``has the customer paid?'' answered without leaving Sales} (d1.north);
\end{tikzpicture}
\end{erpfig}

A strong sales-order system therefore validates master data at entry, applies
configured pricing, checks availability and credit, creates fulfilment documents
from the order, passes billing data to Finance, and records status changes so
Sales can answer customer questions.

\section{Order-to-cash as a cross-functional cycle}

\term{Order-to-cash} emphasises that the process is not complete when the sales
order is confirmed. The business has not achieved what it wanted until the
product is fulfilled and the money is collected.

\begin{erptable}
\begin{tabularx}{\linewidth}{@{}>{\bfseries\leavevmode\color{ERPNavy}}p{0.19\linewidth}p{0.24\linewidth}X@{}}
\toprule
\rowcolor{ERPPaleBlue!92}
Stage & Primary owner & Information created or consumed \\
\midrule
Opportunity / quotation & Sales / Marketing & Customer need, product, quantity, expected price, proposed terms. \\
Sales order & Sales & Confirmed demand and commercial commitment. \\
Availability and delivery & Inventory / Operations & Stock reservation, picking, actual shipped quantity, delivery status. \\
Invoice & Finance & Amount due, tax and payment terms, receivable record. \\
Payment & Finance & Cash received and settlement status, visible back to Sales. \\
\bottomrule
\end{tabularx}
\end{erptable}

\section{Integration of Sales and Accounting}

In an ERP system accounting consequences are created from operational events
rather than reconstructed later. Management reporting, customer balances, and
profitability analysis all depend on that same source transaction.

If a customer receives a discount, Sales and Finance need the same final price.
Re-entering the invoice manually creates an unnecessary opportunity to disagree.

\begin{keyidea}
A reliable sales process preserves document lineage: the quotation explains how
the order was formed, the delivery explains what was physically fulfilled, the
invoice explains what was billed, and the payment explains what was settled.
\end{keyidea}

\section{Document states matter}

Business records move through a lifecycle. A draft is editable; a confirmed or
posted record is a commitment; a done or paid state records completion of the
real-world event. Knowing which state a record is in explains almost every
``the button is missing'' problem in the labs.

\begin{erpfig}[Document states]{The three lifecycles used throughout the
practicals. An action is available only when the record is in the state that
permits it, which is why troubleshooting runs
backwards.}{fig:document-states}
\begin{tikzpicture}[x=1cm,y=1cm]
\foreach \i/\lbl/\a/\b/\c/\ta/\tb in {%
  0/{Sales document}/Quotation/{Sales Order}/{Locked}/{Confirm}/{Fulfilled},
  1/Delivery/Draft/{Ready --- reserved}/Done/{Check availability}/{Validate},
  2/{Customer invoice}/Draft/Posted/{Paid / In Payment}/{Post}/{Pay}}{
  \pgfmathsetmacro{\y}{-\i*1.62}
  \node[erpcap,anchor=west,align=left,text width=2.2cm] at (-0.1,\y) {\bfseries \lbl};
  \node[erpquiet,minimum width=3.0cm,minimum height=0.8cm] (a\i) at (3.7,\y) {\a};
  \node[erpstep,minimum width=3.2cm,minimum height=0.8cm]  (b\i) at (8.5,\y) {\b};
  \node[erpgood,minimum width=3.2cm,minimum height=0.8cm]  (c\i) at (13.1,\y) {\c};
  \draw[erpflow] (a\i)--node[erpnote,above,fill=white,inner sep=1pt]{\ta}(b\i);
  \draw[erpflow] (b\i)--node[erpnote,above,fill=white,inner sep=1pt]{\tb}(c\i);
}
\node[erpcap,anchor=west,color=UTPMuted] at (2.6,1.05) {\bfseries editable};
\node[erpcap,anchor=west,color=UTPBlue]  at (7.4,1.05) {\bfseries committed};
\node[erpcap,anchor=west,color=UTPGreen] at (12.1,1.05) {\bfseries complete};
\end{tikzpicture}
\end{erpfig}

If a button is missing, the first diagnostic question is: \emph{is the preceding
business event complete?}

\section{Company-wide sales analysis}

Tutorial 5 asks how management could analyse the whole relationship with a
customer such as Delicious Foods when each sales division keeps its own data.
In a fragmented environment this requires collecting data from every division,
reconciling customer identities and product names, standardising terms, removing
duplicates, and combining the results --- and divisions may be reluctant to
share performance data at all.

\begin{erpfig}[One customer, three identities]{Why fragmented data resists
consolidation. The three divisions each hold a valid record; only controlled
master data makes them the same customer.}{fig:customer-identity}
\begin{tikzpicture}[x=1cm,y=1cm]
\node[erphead,color=UTPRed] at (3.5,2.6) {Divisional records};
\node[erpbad,minimum width=3.6cm,minimum height=0.72cm] (x1) at (3.5,1.65) {``Delicious Foods''};
\node[erpbad,minimum width=3.6cm,minimum height=0.72cm] (x2) at (3.5,0.7)  {``Delicious Foods Co.''};
\node[erpbad,minimum width=3.6cm,minimum height=0.72cm] (x3) at (3.5,-0.25) {``DELICIOUS FOODS CO''};
\node[erpwarn,minimum width=3.3cm,minimum height=1.05cm] (agg) at (8.4,0.7)
  {Manual roll-up\\\tiny\mdseries counts three customers};
\foreach \n in {x1,x2,x3}{\draw[erpslow] (\n)--(agg);}
\node[erpstore,minimum width=3.7cm,minimum height=1.05cm,draw=UTPGreen,fill=UTPPaleTeal]
  (gov) at (13.3,0.7) {One controlled\\\scriptsize\bfseries customer identity};
\draw[erpflow] (agg)--node[erpnote,above,text width=1.9cm]{master-data governance}(gov);
\node[erpnote,text width=4.0cm] at (13.3,-0.7)
  {transactions reference the identity, so the roll-up is automatic};
\end{tikzpicture}
\end{erpfig}

ERP reduces the technical burden by storing comparable data centrally. It does
not remove organisational reluctance, but it makes a unified view possible where
governance and management expectations support one.

\section{How the sequence appears in Odoo 19 Community}

\begin{odoobox}
For the Road Bike example, the practical sequence is:
\begin{enumerate}
  \item create a CRM opportunity for Seri Iskandar Cycling Club;
  \item create a quotation for 5 Road Bikes at RM 1,000 each;
  \item confirm the quotation so it becomes a Sales Order;
  \item open the linked Delivery and validate shipment of 5 bikes;
  \item return to the Sales Order and create a customer invoice;
  \item post the invoice;
  \item register the payment so the invoice becomes Paid or In Payment.
\end{enumerate}
The price is set in Sales and flows into the invoice rather than being typed
again in Finance.
\end{odoobox}

\begin{tryit}
Take the old Fitter Snacker process and label every point where information is
copied, delayed, estimated, or re-entered. Map each problem to one ERP
mechanism: shared master data, automatic document creation, real-time status,
configured pricing, availability checking, or integrated accounting.
\end{tryit}

\begin{quickcheck}
\begin{enumerate}
  \item Why can a quotation error later become an invoicing problem?
  \item What is the difference between a Sales Order and a Delivery in the process model?
  \item Why is payment status useful to Sales even though Finance records the payment?
\end{enumerate}
\end{quickcheck}

\begin{chapterrecap}
\begin{itemize}
  \item The sales order process spans quotation, order, fulfilment, billing, payment, and returns, with a different owner at almost every step (\figref{fig:order-to-cash-lanes}).
  \item Fitter Snacker's problems all come from one cause: separate copies of the same order (\figref{fig:fitter-snacker-silos}).
  \item Most of an unintegrated order's elapsed time is waiting, not work (\figref{fig:fitter-timeline}).
  \item ERP value comes from lineage --- each document inherits its data from the one before (\figref{fig:document-lineage}).
  \item Document states decide which actions are available, which makes troubleshooting a backwards walk (\figref{fig:document-states}).
\end{itemize}
\end{chapterrecap}
